% !TEX root = ../samplepaper.tex
% Methodology section for the MACLEAN'26 paper (LNCS / llncs).
% Designed to be \input{} into samplepaper.tex. Assumes the preamble already
% loads: amsmath, amssymb, graphicx, booktabs. Citation keys are placeholders;
% reconcile them with your .bib (he2022mae, snell2017proto, lin2023ssmae,
% mft, hypersigma, houston2013, trento, muufl).
%
% NOTE: no architecture figure exists in the repo yet. A placeholder is included
% at the end of Sec. \ref{subsec:encoder}; replace it with a real diagram.

\section{Methodology}
\label{sec:method}

We study few-shot, multimodal land-cover classification and compare two
strategies for obtaining a transferable encoder: a compact, purpose-built
spectral--spatial fusion transformer pretrained with self-supervised masked
modeling (the \emph{CoFFE encoder}), and a large pretrained HSI foundation model
(HyperSIGMA) lightly adapted to each scene. Both are evaluated under an identical
parameter-free episodic protocol, isolating representation quality from
task-specific adaptation.

\subsection{Problem setup and notation}
\label{subsec:setup}
A labelled example is a spatial patch centred on a pixel,
$x = \{x^{\text{hsi}}\!\in\!\mathbb{R}^{C_h\times P\times P},\,
        x^{\text{aux}}\!\in\!\mathbb{R}^{C_a\times P\times P}\}$,
with patch size $P{=}11$, $C_h$ hyperspectral bands and $C_a$ auxiliary (LiDAR
elevation) channels. Evaluation is episodic: each $N$-way $K$-shot episode draws
a support set $S$ of $K{=}5$ labelled patches per class and a query set $Q$. A
frozen encoder $f_\theta$ maps a patch to a feature; the class prototype is the
mean of its support features,
\begin{equation}
  c_n = \frac{1}{K}\!\!\sum_{(x_i,y_i)\in S,\, y_i=n}\!\! f_\theta(x_i),
  \label{eq:prototype}
\end{equation}
and a query is assigned to the nearest prototype~\cite{snell2017proto}. There is
\textbf{no gradient finetuning at test time} and the classifier has no learnable
parameters, so the only transferred knowledge is the pretrained encoder. We
report Overall Accuracy (OA), Average per-class Accuracy (AA), and Cohen's
$\kappa$ as means over episodes with $95\%$ confidence intervals; query sets are
class-balanced, hence OA${\equiv}$AA and we report OA and~$\kappa$.

\subsection{Datasets and preprocessing}
\label{subsec:data}
We use Houston~2013 ($C_h{=}144$, $C_a{=}1$, $15$ classes), Trento ($63,1,6$) and
MUUFL ($64,2,11$). Each scene is cut into $11\times11$ patches centred on every
pixel; every HSI band and auxiliary channel is min--max normalised to $[0,1]$.
Self-supervised pretraining uses all pixels (labelled and unlabelled).

\subsection{The CoFFE encoder: unified multimodal fusion}
\label{subsec:encoder}
The encoder performs \emph{early, token-level fusion}. HSI and auxiliary bands are
concatenated along the channel axis and embedded jointly: a $1{\times}1$
convolution (BatchNorm, GELU) projects the $C_h{+}C_a$ bands to dimension $D$, and
a $3{\times}3$ convolution (BatchNorm, GELU) mixes local context before flattening
to one token per pixel,
\begin{equation}
  [B, C_h{+}C_a, P, P] \;\rightarrow\; [B, D, P, P] \;\rightarrow\;
  [B, N{=}P^2, D].
\end{equation}
Each of the $N{=}121$ tokens thus carries both modalities at its pixel. A learnable
class-agnostic token is prepended, learnable positional embeddings are added, and
the sequence is processed by a small pre-LayerNorm transformer (multi-head
self-attention and a $4\times$ GELU MLP per block). The labelled budget is tiny, so
we use a deliberately compact configuration ($D{=}128$, $2$ heads, $2$ layers). A
projection-head MLP shapes the pretraining objective but is \textbf{discarded for
evaluation}; few-shot features are taken from the encoder output. We do not use any
class-aware patch-embedding adaptation.

% TODO: replace with a real architecture diagram.
% \begin{figure}[t]\centering
%   \includegraphics[width=\linewidth]{figures/coffe_architecture.pdf}
%   \caption{The CoFFE encoder and the adapted HyperSIGMA baseline.}
%   \label{fig:arch}
% \end{figure}

\subsection{Self-supervised pretraining}
\label{subsec:pretrain}
The encoder is pretrained by masked reconstruction over the unified token space,
turning unlabelled pixels into supervision. Two composable masks act on the
concatenated input: \textbf{band masking} hides each $(\text{pixel},\text{band})$
entry independently with probability $r_b$ (masked entries replaced by a learnable
per-channel fill), and \textbf{spatial-token masking} replaces a fraction $r_s$ of
pixel tokens with a learnable mask token after tokenization. A lightweight $2$-layer
MLP decoder reconstructs the full $C_h{+}C_a$-dimensional pixel vector, and the loss
is a masked mean-squared error over the union of the two masks,
\begin{equation}
  \mathcal{L} \;=\; \frac{\sum_j w_j\, m_j\, (\hat v_j - v_j)^2}{\sum_j m_j},
  \label{eq:recon}
\end{equation}
with optional Gaussian center-weighting $w_j$ that emphasises the patch centre (the
classified pixel). We study three regimes: \emph{spectral} $(r_b,r_s){=}(0.85,0)$,
\emph{spatial} $(0,0.75)$, and \emph{combined} $(0.85,0.75)$. As a strong
single-objective baseline we also pretrain a standard MAE~\cite{he2022mae} in the
same token space ($75\%$ token dropping, asymmetric transformer decoder). All
encoders use AdamW (lr $1.5{\times}10^{-4}$, cosine schedule, $100$-epoch warm-up,
weight decay $0.05$, gradient clip $1.0$, batch $64$; $3000$ epochs on Houston,
$1500$ on Trento/MUUFL).

\subsection{Parameter-free few-shot evaluation}
\label{subsec:eval}
At test time the encoder is frozen and the projection head removed. Patch tokens are
pooled (uniform or Gaussian center-weighted) to one vector per patch; prototypes
follow Eq.~\eqref{eq:prototype}. We use a Euclidean prototype rule as the primary
classifier, $\text{logit}_n = -\lVert z_q - c_n\rVert_2^2$, and report a cosine
variant ($\ell_2$-normalised features, $\text{logit}_n = \tau\,\langle z_q,c_n\rangle$,
$\tau{=}10$) for comparison; Euclidean is consistently stronger on these features.
Episodes are $5$-shot, $N$-way ($N$ = number of classes), with $100$ queries per
class over $1000$ episodes.

\subsection{Foundation-model baseline: adapted HyperSIGMA}
\label{subsec:hypersigma}
As the foundation-model point of comparison we use HyperSIGMA~\cite{hypersigma}, a
dual-branch ViT. A \emph{spatial} branch encodes a $3$-component PCA of the HSI with
a frozen ViT-Base body; a \emph{spectral} branch pools the raw bands to $100$
spectral tokens and encodes them with a second frozen ViT-Base; a gated
spatial--spectral enhancement module (SEM) fuses the branches into a $512$-d
feature. Only the randomly-initialised input/projection layers and the SEM/decoders
are trainable; the ViT bodies remain frozen. We adapt these trainable pieces to each
scene's unlabelled pixels by a light MAE finetuning ($75\%$ token masking,
per-token normalised targets), with three modes---\texttt{spatial\_only},
\texttt{spectral\_only}, and \texttt{joint\_sem} (both branches plus SEM)---and
evaluate the corresponding feature (\texttt{fused}, \texttt{spat\_pool}, or
\texttt{spec\_pool}) under the identical protocol of
Section~\ref{subsec:eval}.
